% ============================================================================
% 8 DISCUSSION
% ============================================================================
\section{Discussion}
\label{sec:discussion}

\paragraph{RQ1: reliable output is not the same as valid output.}
The short-input baselines demonstrate two different ceilings. Single-type
Few-NERD extraction reaches about $0.79$ strict $F_1$, while synthetic-ticket
routing reaches only $0.422$ accuracy despite perfect JSON and allowed-label
validity. The agentic classifier adds a call but does not improve accuracy. The
proper conclusion is narrow: on this model, prompt and dataset, extra intrinsic
reasoning is not evidence of better predictions, consistent with findings on
self-correction without external feedback \citep{huang2024selfcorrect}. The
unclean 1{,}129-tag vocabulary and synthetic origin of the ticket corpus also
mean this is a stress test, not an estimate for an internal ticket system.

\paragraph{RQ2: useful pre-annotation, with a substantial review burden.}
The fresh SciREX confirmation gives the cleanest quality estimate: $0.5738$
strict and $0.7030$ overlap $F_1$. Thus the model often points to a relevant
phrase without reproducing the annotation boundary exactly. That distinction
matches the product goal. A reviewer can benefit from useful localization, but
the output is not yet suitable as unattended training data. Metric and Dataset
remain the hardest labels, and the 100-paper run shows Dataset precision of only
$0.215$ despite recall of $0.642$. Better negative definitions and reviewed
in-context examples should target this error before changing models.

The length analysis is encouraging operationally. All full windows, including
at least 201 sentences in the very-long bucket, completed and overlap $F_1$ was
stable. It does not show that more document context helps because each sentence
is predicted independently. Rather, it shows that length no longer causes silent
truncation or unrecoverable execution. Cross-sentence coreference and relation
extraction remain separate problems, and our mention score is not directly
comparable with SciREX's original document-level relation pipeline
\citep{jain2020scirex}.

\paragraph{RQ3: a reusable path to gold data is implemented, not yet validated
as a productivity gain.}
The product now covers configuration, prompt inspection, pre-annotation,
LLM or local-embedding classification suggestions, human correction,
provenance and export for more than one domain. It also makes a crucial semantic
distinction: uploaded raw text
contains no hidden gold labels. Gold comes either from an imported annotated
dataset or from reviewer-approved corrections. Department routing is therefore
stored separately as model prediction and approved decision.

This proves functional completeness, not annotation-time savings. The largest
remaining empirical gap is the human-in-the-loop workflow itself. A valid study
should use unseen, representative internal tickets in a randomized,
counterbalanced comparison of manual annotation and pre-annotation review.
At least two domain experts should independently annotate an overlapping subset,
with disagreements adjudicated under a frozen guideline. The application should
record active review time excluding idle periods, additions, deletions, boundary
changes and relabeling. Primary outcomes should include time per document and
final strict quality against adjudicated labels; agreement should be reported
with an appropriate chance-corrected statistic and pairwise span scores
\citep{artstein2008intercoder}. A pilot should inform power and sample-size
planning. Acceptance rate alone is insufficient because a reviewer can accept
wrong but plausible suggestions.

Classification is less validated than extraction. The 1{,}000-ticket comparison
uses synthetic data, while the local embedding route and retrieval-augmented LLM
route have no independent target-domain result. A target-ticket study should
freeze the queue schema and partition data by customer, conversation thread or
time into non-overlapping reference, development and held-out test sets. The
reference set may populate the embedding index, the development set may select
the prompt and retrieval settings, and the held-out set should be scored once.
LLM and embedding routes should be compared on the same tickets using accuracy,
macro-$F_1$, per-queue confusion, abstention or low-confidence coverage, latency
and token use. Gold queue labels must be assigned before reviewers see either
prediction.

\begin{takeaway}[title={Engineering recommendations}]
\small
\textbf{1. Begin with an explicit schema and a small, independent gold set.}
Use label definitions, positive and negative examples, and exact boundary rules;
do not treat well-formed JSON as evidence of correctness.

\textbf{2. Report strict and reviewer-oriented metrics together.} Strict spans
measure training-data readiness, while boundary-tolerant and overlap scores
measure whether pre-annotation points a human to the right phrase.

\textbf{3. Treat corrections as data.} Hold out part of the reviewed corpus for
evaluation, retrieve approved examples for prompting, and consider fine-tuning
only after the annotation policy and inter-annotator agreement are stable.

\textbf{4. Make long runs restartable and cost-visible.} Checkpoint per sentence,
hash the prompt and manifest, record token use, and configure actual provider
prices before presenting cost. A stored zero with zero rates means unknown, not
free.
\end{takeaway}

\subsection{Limitations and threats to validity}

\textbf{External validity.} Few-NERD is Wikipedia, SciREX is scientific prose,
and the ticket corpus is synthetic. None establishes accuracy on an
organization's internal tickets or service reports. All LLM results use one
hosted deployment, so model rankings cannot be generalized. SciREX's released
mentions combine automatic and human annotation and can contain guideline or
annotation noise.

\textbf{Internal and statistical validity.} Most configurations were run once
at temperature $0$; deterministic decoding reduces but does not eliminate
provider-side variation. The 20-paper held-out interval is wide, and the
development prompt comparison used only ten detected sentences per window.
The 100-paper operational sample is drawn from the training split and is a
reliability and scale gate, not a second held-out test. The 1{,}000-window corpus
contains overlapping windows from 438 papers, so window-level independence
would be false; this is why intervals cluster by \code{doc\_id}. The full
1{,}000-window inference run is prepared but incomplete.

\textbf{Metric validity.} Boundary-tolerant and overlap matching are deliberately
more forgiving and should never replace strict scoring in claims about gold-data
quality. Conversely, strict $F_1$ understates pre-annotation utility when the
reviewer can repair a nearby boundary quickly. Neither measure evaluates time
saved, cognitive load, or missed cross-sentence relations.

\textbf{Product scope.} The application is a reproducible single-instance
research product. It lacks authentication, role-based permissions, concurrent
work allocation, a central database, tenant isolation, and direct integration
with a ticketing platform. Reviewer corrections are logged but do not
automatically retrain a model. Those controls are required before deployment to
hundreds of users or processing confidential tickets.
